\section{Method}
\label{sec:method}

\subsection{Problem Formulation}
\label{sec:problem}

We consider a cooperative interception task with three blue UAV actors---scout, relay, and attacker---and one target. The target is represented as a graph node but is not an actor. The problem is formalized as a decentralized partially observable Markov decision process (Dec-POMDP) $\langle \mathcal{I}, \mathcal{S}, \mathcal{A}, O, P, R, \gamma \rangle$ with agents $\mathcal{I}=\{0,1,2\}$ (scout, relay, attacker). At each step, agent $i$ obtains a local observation $o_i$ that includes an ego-state slice, an implicit role identity (role one-hot), and a local observation of the multi-relation graph; the team takes joint actions under a shared cooperative reward that encourages closing the kill chain, avoiding collisions, and maintaining task-relevant information flow.

The environment includes two sources of degradation relevant to this paper. First, communication is unreliable: messages may be dropped (dropout) or arrive with delay, and the physical communication availability is represented by an environment-driven adjacency $\mathrm{comm\_adj}$ rather than by a learned gate. Second, a blue node (typically the relay) can fail for a temporary window, during which it neither sends nor relays information. The optimization target is therefore not only to maximize the terminal success probability but also to minimize the post-failure coordination-recovery latency: after a node or link failure, the team should re-establish effective coordination as quickly as possible. We define failure exposure and recovery conceptually here; the statistical definitions used for evaluation are given in Section~\ref{sec:setup}.

\subsection{Task-Graph Construction}
\label{sec:task_graph}

The graph has $N=4$ nodes: the three blue actors and the target. A multi-relation adjacency $\mathrm{relation\_adj} \in \mathbb{R}^{3\times N\times N}$ encodes three relation channels, one per relation $r \in \mathcal{R}=\{\mathrm{per}, \mathrm{comm}, \mathrm{sup}\}$ (perception, communication, task support), with the convention $A_{ij}=1$ meaning that sender $j$ connects to receiver $i$.

\begin{itemize}
    \item \textbf{Perception}: direct sensing edges, active when the sender can detect or track the target;
    \item \textbf{Communication}: edges gated by the environment-driven communication availability $\mathrm{comm\_adj}$ (a function of distance, dropout realization, and failure state), i.e., whether a message is physically deliverable;
    \item \textbf{Task support}: task-dependency edges (see Section~\ref{sec:task_support}).
\end{itemize}

The task graph represents task dependency structure; it is not itself the physical communication network. In particular, $\mathrm{comm\_adj}$ is provided by the environment and is not a learned quantity. The union of the relation channels (together with a self-loop) forms the node mask used during message aggregation.

\subsection{Multi-Relation Graph Representation}
\label{sec:graph_rep}

Each relation channel is processed by an independent role-conditioned graph-attention layer, and the encoder stacks two such layers, each containing the three relation channels. Each layer also contains a separate union-graph attention branch whose output is fused with the relation outputs as a global residual information path. For receiver node $i$ and sender node $j$, relation $r$ has adjacency $A^r_{ij}$, where $A^r_{ij}=1$ means that sender $j$ is available to receiver $i$. The projected sender representation is $h'_j=W h_j$; edge features do not concatenate to the message payload. Instead, they modulate the attention score, while the role-pair gate modulates the sender message. The relation message is
\begin{equation}
    m^r_{ij} = A^r_{ij}\,\alpha^r_{ij}\,\mu^r_{ij}\,h'_j,
\label{eq:message}
\end{equation}
where $\alpha^r_{ij}$ is the relation-specific attention weight and $\mu^r_{ij}$ is the learned receiver--sender role-pair modulation. The receiver update aggregates the relation-specific outputs and the union-branch output as
\begin{equation}
    \tilde{h}_i = \phi\left(h_i, \operatorname{Fuse}\left(\left\{\sum_j m^r_{ij}\right\}_{r\in\mathcal{R}},\;m^{\mathrm{union}}_i\right)\right),
\label{eq:aggregate}
\end{equation}
where $\phi$ is a non-linear update with a residual addition of the input node representation. The single-graph baseline replaces the relation-specific encoder and union branch with ordinary graph-attention layers over the merged adjacency. The no-graph ablation disables graph aggregation within the RI-MAPPO actor pipeline, whereas canonical MAPPO is an independently configured no-graph CTDE baseline.

\begin{figure}[t]
    \centering
    \includegraphics[width=0.98\linewidth]{intercept_3d_multi_relation_graph.png}
    \caption{Multi-relation task graph. Perception, communication, and task-support relations have different semantics and are processed by state-dependent, edge-feature-modulated attention. Task-support is a task-dependent relational mask over delivered communication rather than an independent information channel; a static role-pair modulation is applied as an auxiliary component.}
    \label{fig:multi_relation_graph}
\end{figure}

\subsection{State-Dependent Edge-Feature-Modulated Attention}
\label{sec:attention}

The attention weight in Eq.~\eqref{eq:aggregate} is state-dependent and edge-feature-modulated. The implementation uses one shared node projection $W$ within each attention layer rather than independent query and key projections. For relation $r$, the raw attention score between receiver $i$ and sender $j$ is
\begin{equation}
    s^r_{ij} = \mathrm{LeakyReLU}\big(a^r \cdot [W h_i \,\|\, W h_j]\big) + b(e_{ij}),
\label{eq:score}
\end{equation}
where $e_{ij}$ is the edge feature vector (relative position, velocity, altitude, bearing, communication reachability, message age, information confidence, sensing source, and kill-chain stage) and $b(\cdot)$ is a learned scalar edge-score network. The graph edge-feature dimension ($17$) describes \emph{graph computation} only; it is not the communication payload. The relation adjacency acts as a hard mask before the softmax: scores of disconnected pairs are set to $-\infty$, and
\begin{equation}
    \alpha^r_{ij} = \mathrm{softmax}_j\big( s^r_{ij} + \log \hat{A}^r_{ij}\big), \qquad \hat{A}^r_{ij}=\begin{cases}1,& A^r_{ij}=1 \text{ or } i=j,\\0,&\text{otherwise.}\end{cases}
\label{eq:softmax}
\end{equation}
Physical message delivery is governed separately by the environment's $\mathrm{comm\_adj}$; the attention mask reflects graph-relation availability and may include the task-support channel, but it does not itself decide which bytes are transmitted. Consequently, dynamic adaptation in EA-RG arises primarily from state-dependent graph attention, relation masks, edge-score biases, and environment-driven communication availability; the role-pair modulation is static with respect to the input state and is not a communication-pruning mechanism.

\subsection{Task-Support Relation}
\label{sec:task_support}

The task-support channel is the slice $\mathrm{relation\_adj}[:,2,:,:]$, with receiver--sender orientation. An edge is active when (i) the receiver--sender pair is role-compatible with a task-dependency (e.g., scout or relay holding target information toward an attacker), (ii) the corresponding communication edge is available, and (iii) the sender currently holds task-relevant information (target information, or a local attack window for an attacker). Task-support therefore biases relational reasoning toward task-relevant delivered information; it is a task-dependent relational mask over delivered communication and is not an independent communication channel, nor does it decide whether real messages are sent. We do not claim any specific post-failure activation behavior for this channel here; its empirical contribution is evaluated in Section~\ref{sec:setup} and its temporal analysis in the Discussion.

\subsection{Gate Prior and Static Role-Pair Modulation}
\label{sec:gate_prior}

This section separates two concepts that are often conflated.

\textbf{Static role-pair modulation.} Each relation channel additionally multiplies the sender message by a learned gate
\begin{equation}
    \mu^{r}_{ij} = \sigma\big( \mathrm{Emb}^{(r)}_{\rho_i,\rho_j} \big),
\label{eq:rpg}
\end{equation}
where $\rho_i, \rho_j$ are receiver and sender roles, $\mathrm{Emb}^{(r)}$ is a learned role-pair embedding for relation $r$, and $\sigma$ is the logistic function. This is a \emph{learned static role-pair modulation}: the embedding is a function of the role pair and relation only, not of the failure state or of any input feature. It performs multiplicative modulation, does not prune messages, and does not change the number of real transmitted messages. The modulation itself is retained as an auxiliary architectural component; its independent empirical contribution is evaluated separately.

\textbf{Gate Prior.} The gate logits are initialized to a fixed structured value before training. In the full method a subset of role-pair gate logits is set to $0.4$ (sigmoid $\approx 0.599$), while in the w/o Gate Prior variant all logits are initialized to $0$ (sigmoid $= 0.5$). The prior therefore acts as a structured initialization / optimization bias over static role-pair modulation; it is not a runtime mechanism. We defer the question of whether this prior empirically improves stability to RQ6 (Section~\ref{sec:experiments}).

\subsection{Actor--Critic Optimization}
\label{sec:optim}

The policy is trained with a MAPPO-style clipped objective under centralized training and decentralized execution:
\begin{equation}
    \mathcal{L}_{\mathrm{PPO}} =
    \mathbb{E}\left[
        \min\left(
            r_t(\theta)\hat{A}_t,
            \mathrm{clip}(r_t(\theta),1-\epsilon,1+\epsilon)\hat{A}_t
        \right)
    \right],
\label{eq:ppo}
\end{equation}
where $r_t(\theta)$ is the policy probability ratio and $\hat{A}_t$ is the advantage estimate. The \emph{actor} receives the local observation (including the role one-hot), the node features, edge features, role vector, the union adjacency, and the multi-relation adjacency, and outputs a joint action in a decentralized manner. The \emph{critic} is a centralized value function that takes the global shared observation concatenated with the role one-hot; it does not use the graph gate. Behavior-cloning warm-start from the pretrained teacher is used as a training aid in the formal method chain. Detailed training schedules (number of seeds and updates) are given in Section~\ref{sec:setup}.
